\documentclass{article}
\usepackage{amsmath}
\begin{document}
\section{线性回归}
线性回归是最简单的回归模型之一，它假设输入特征与输出之间存在线性关系。线性回归的目标是找到一组权重，使得预测值与实际值之间的误差最小化。
\subsection{模型定义}
线性回归模型可以表示为：
\begin{equation}
y = w^T x + b
\end{equation}
其中，$y$ 是预测值，$x$ 是输入特征向量，$w$ 是权重向量，$b$ 是偏置项。
\subsection{损失函数}
线性回归通常使用均方误差（Mean Squared Error, MSE）作为损失函数，定义如下：
\begin{equation}L(w, b) = \frac{1}{N} \sum_{i=1}^{N} (y_i - (w^T x_i + b))^2\end{equation}
其中，$N$ 是样本数量，$y_i$ 是第 $i$ 个样本的实际值，$x_i$ 是第 $i$ 个样本的输入特征向量。
\subsection{参数优化}
为了最小化损失函数，我们可以使用梯度下降算法来更新权重和偏置项。更新规则如下：
\begin{equation}w := w - \alpha \frac{\partial L}{\partial w}\end{equation}
\begin{equation}b := b - \alpha \frac{\partial L}{\partial b}\end{equation}
其中，$\alpha$ 是学习率，控制更新的步长。
\subsection{正则化}
为了防止过拟合，线性回归模型可以添加正则化项。常见的正则化方法有L1正则化（Lasso回归）和L2正则化（Ridge回归）。L1正则化的损失函数定义为：
\begin{equation}L(w, b) = \frac{1}{N} \sum_{i=1}^{N} (y_i - (w^T x_i + b))^2 + \lambda \sum_{j=1}^{d} |w_j|\end{equation}
L2正则化的损失函数定义为：
\begin{equation}L(w, b) = \frac{1}{N} \sum_{i=1}^{N} (y_i - (w^T x_i + b))^2 + \lambda \sum_{j=1}^{d} w_j^2\end{equation}
其中，$\lambda$ 是正则化参数，控制正则化项的权重。
\subsection{模型评估}线性回归模型的评估指标包括均方误差（MSE）、平均绝对误差（MAE）和决定系数（R-squared）。这些指标可以帮助我们评估模型的性能和预测能力。均方误差越小，模型的预测能力越好；平均绝对误差越小，模型的预测误差越小；决定系数越接近1，模型的拟合程度越好。
\end{document}